% Section 9.3, from lectures/L09/appendix/b_exercises.md.

\section{Exercises}
\label{c9:sec:exercises}\label{c9:app:b}

\begin{aside}
\textbf{How to work these.} The Reflection exercises are answered on paper, from the chapter. The
Code exercises are programs for the board: build and flash them the way \secref{c2:app:c}
(Microchip Studio) or \secref{c2:app:d} (Linux) describes, or run them in the simulator of
Appendix~\ref{sim:ch} if you have no board yet.

\textbf{How to check your work.} The Reflection exercises, and the questions Code
Exercises~\ref{c9:ex:3} to~\ref{c9:ex:6} ask about their design, are answered in
Appendix~\ref{app:answers}. The Code exercises have no printed solution: the board tells you
whether a program does what the exercise describes. Try each exercise before you read its answer.
\end{aside}

% ------------------------------------------------------------------------------------------
\exerciseset{Concepts}
\label{c9:set:1}

\reflexercise{A struct or a macro}
\label{c9:ex:1}
What problem does bundling a pin number into a \code{gpio_t} struct solve that a
\code{#define}d \code{LED_ON} macro can't? What has to change to support a second LED under each
approach?

\reflexercise{By pointer or by value}
\label{c9:ex:2}
\code{gpio_read()} takes \code{const gpio_t* self} rather than \code{gpio_t self} (the whole
struct, by value). What would change, functionally or in cost, if it took \code{gpio_t self}
instead? Would \code{gpio_init()} still work correctly if it took \code{self} by value?

% ------------------------------------------------------------------------------------------
\exerciseset{Program}
\label{c9:set:2}

\codeexercise{Freeing a GPIO}
\label{c9:ex:3}
Add \code{gpio_free(gpio_t* self)} to \secref{c9:a:2}'s driver:
\begin{itemize}
  \item Reset the pin's \code{DDRx}/\code{PORTx} bits back to their power-on state (input,
    pull-up disabled).
  \item Decide whether \code{self->ddrx}, \code{self->portx}, and \code{self->pinx} should be set
    to \code{NULL} afterwards.
  \item Justify the choice:
  \begin{itemize}
    \item What does calling \code{gpio_write()} or \code{gpio_toggle()} on a freed instance do in
      each case?
    \item Which is preferable?
  \end{itemize}
\end{itemize}

\codeexercise{Initializing from a pin number}
\label{c9:ex:4}
Rewrite \code{gpio_init()} so callers pass a single pin number, 0--13, instead of
\code{DDRx}/\code{PORTx}/\code{PINx} pointers:
\begin{itemize}
  \item Pin mapping:
  \begin{itemize}
    \item 0--7 map to \code{PORTD0}--\code{PORTD7}.
    \item 8--13 map to \code{PORTB0}--\code{PORTB5}.
  \end{itemize}
  \item \code{gpio_init()} should compute and store \code{ddrx}/\code{portx}/\code{pinx} (and the
    bit position) itself.
  \item Return \code{bool}:
  \begin{itemize}
    \item \code{false} for any pin number outside 0--13, without touching any register.
    \item \code{true} otherwise, after initialization.
  \end{itemize}
  \item \code{pin} is \code{const} in \secref{c9:a:2}'s \code{gpio_t} because it's only ever set
    at struct-literal time; now that \code{gpio_init()} derives it from the pin number instead,
    does it still make sense for \code{pin} to stay \code{const}?
\end{itemize}

\codeexercise{A pin-reservation table}
\label{c9:ex:5}
Add a pin-reservation table:
\begin{itemize}
  \item Declare a \code{static uint16_t pin_reg} in \file{gpio.c}.
  \item Use one bit for each pin number introduced in Exercise~\ref{c9:ex:4}:
  \begin{itemize}
    \item Bit~$n$ represents pin~$n$.
    \item \code{1} means that the pin is reserved.
    \item \code{0} means that the pin is free.
  \end{itemize}
  \item Update \code{gpio_init()}:
  \begin{itemize}
    \item Check the bit corresponding to the requested pin.
    \item If the bit is already set, return \code{false}.
    \item In that case, do not modify any hardware register.
    \item If the bit is clear, set it to reserve the pin.
    \item Then initialize the GPIO as in Exercise~\ref{c9:ex:4}.
    \item Return \code{true} when the initialization succeeds.
  \end{itemize}
  \item Extend the \code{gpio_free()} function from Exercise~\ref{c9:ex:3}:
  \begin{itemize}
    \item Clear the pin's bit in \code{pin_reg}.
    \item This allows a later call to \code{gpio_init()} to reserve the same pin again.
  \end{itemize}
  \item Answer the following questions:
  \begin{itemize}
    \item Why must \code{pin_reg} be declared \code{static}?
    \item What would break if it were not \code{static}?
  \end{itemize}
\end{itemize}

\codeexercise{A timer driver}
\label{c9:ex:6}
Capstone --- harder: Implement a timer driver named \code{timer_t}:
\begin{itemize}
  \item Follow the header/source separation used in \secref{c9:a:3}:
  \begin{itemize}
    \item \file{driver/timer.h}
    \item \file{driver/timer.c}
  \end{itemize}
  \item The timer must run a callback function whenever its timeout expires.
  \item \code{timer_t} should follow the same ownership model as \code{gpio_t}:
  \begin{itemize}
    \item The caller declares and owns each \code{timer_t} instance.
    \item The caller may declare an array of timer instances, just as \secref{c9:a:4} declares an
      array of \code{gpio_t} instances.
    \item \code{timer_init()} initializes an existing instance in place.
  \end{itemize}
  \item The ATmega328P has only three hardware timer/counters:
  \begin{itemize}
    \item At most three \code{timer_t} instances may therefore be active at the same time.
    \item Use a reservation scheme similar to Exercise~\ref{c9:ex:5}'s \code{pin_reg}.
    \item \code{timer_init()} must not create more active timer instances than there are hardware
      timer/counters available.
  \end{itemize}
  \item \code{timer_t} is a complete timer driver instance:
  \begin{itemize}
    \item Each instance has:
    \begin{itemize}
      \item Its own callback.
      \item Its own elapsed-time count.
      \item Its own start/stop state.
      \item One reserved hardware timer/counter.
    \end{itemize}
    \item It is driven directly by the hardware timer/counter and its compare-match interrupt.
  \end{itemize}
\end{itemize}

\task{Recommended \file{driver/timer.h}}
\code{timer_callback_t} is a function pointer, see \secref{c1:a:8} if that syntax is unfamiliar.
Document every declaration with Doxygen comments, the same way
\secref{c9:a:2}/\secref{c9:a:3}'s \file{gpio.h} does, and add the header guard yourself (suggested
guard: \code{DRIVER_TIMER_H_}).

\begin{ccode}
#include <stdbool.h>
#include <stdint.h>

/** Alias for timer callbacks. */
typedef void (*timer_callback_t)(void);

typedef struct
{
    // Left for you to design.
} timer_t;

bool timer_init(timer_t* self, uint32_t timeout_ms, timer_callback_t callback);
void timer_clear(timer_t* self);
bool timer_running(const timer_t* self);
void timer_start(timer_t* self);
void timer_stop(timer_t* self);
void timer_toggle(timer_t* self);
void timer_reset(timer_t* self);
\end{ccode}

\task{Initialization and reservation}
\begin{itemize}
  \item \code{timer_init()} must return \code{false} when:
  \begin{itemize}
    \item \code{timeout_ms == 0}.
    \item \code{callback == NULL}.
    \item No hardware timer/counter is available.
  \end{itemize}
  \item Otherwise, \code{timer_init()} must:
  \begin{itemize}
    \item Reserve one free hardware timer/counter.
    \item Associate it with \code{self}.
    \item Store the timeout and callback.
    \item Initialize the timer instance.
    \item Return \code{true}.
  \end{itemize}
  \item \code{timer_clear()} must:
  \begin{itemize}
    \item Stop the timer.
    \item Release its reserved hardware timer/counter.
    \item Allow that hardware timer/counter to be claimed by a later \code{timer_init()} call.
    \item Work whether the later call uses the same \code{timer_t} instance or another one.
  \end{itemize}
  \item Consider the same design question raised for \code{gpio_free()} in
    Exercise~\ref{c9:ex:3}:
  \begin{itemize}
    \item Should \code{timer_clear()} also reset the fields of \code{self}?
    \item If so, which values would leave the instance in a safe, inert state?
  \end{itemize}
\end{itemize}

\task{Timer control}
\begin{itemize}
  \item \code{timer_start()} must start or enable the timer.
  \begin{itemize}
    \item It must not change the elapsed-time count.
  \end{itemize}
  \item \code{timer_stop()} must stop or disable the timer.
  \begin{itemize}
    \item It must not change the elapsed-time count.
  \end{itemize}
  \item \code{timer_toggle()} must switch between the running and stopped states.
  \begin{itemize}
    \item It must not change the elapsed-time count.
  \end{itemize}
  \item \code{timer_reset()} must:
  \begin{itemize}
    \item Set the elapsed-time count to zero.
    \item Preserve the timer's current running or stopped state.
  \end{itemize}
\end{itemize}

\task{Driver independence}
\begin{itemize}
  \item The timer driver must not know that \code{gpio_t} exists.
  \item \file{timer.c} must not directly control an LED or read a button.
  \item Connecting a timer callback to an LED is the responsibility of \file{main.c}.
\end{itemize}

\task{Hardware demonstration}
Wire the following circuit:
\begin{itemize}
  \item Connect a button to D13.
  \item Connect an LED to D8.
  \item Configure a timer with a timeout of 100\,ms.
\end{itemize}
The program must behave as follows:
\begin{itemize}
  \item Pressing the button toggles the timer between running and stopped.
  \item While the timer is running:
  \begin{itemize}
    \item Its callback runs every 100\,ms.
    \item The callback toggles the LED.
  \end{itemize}
  \item While the timer is stopped:
  \begin{itemize}
    \item The LED remains off.
    \item Calling the callback manually must have no effect.
  \end{itemize}
  \item To enforce this behavior, the callback itself must check \code{timer_running()} before
    toggling the LED.
\end{itemize}

\task{Testing}
Confirm \code{timer_init()} returns \code{false} for each of the three invalid cases in
``Initialization and reservation'': a \code{0} timeout, a \code{NULL} callback, and a fourth
\code{timer_t} once three hardware timer/counters are already reserved. Then confirm
\code{timer_clear()} actually frees a hardware timer/counter for reuse: after reserving all three,
clear one and check that a new \code{timer_init()} call succeeds again.
